% -*- ispell-dictionary: "french" -*-

\documentclass{beamer}
\usetheme{metropolis}

\usepackage[french]{babel}
\usepackage[utf8]{inputenc}
\usepackage[T1]{fontenc}

% Listes sur 2 colonnes
\usepackage{multicol}

% Tableaux
\usepackage{array}
\usepackage{multirow}

% Images, couleurs
\usepackage{graphicx}
\usepackage{xcolor}

% Mathématiques
\usepackage{mathtools,amssymb,stmaryrd}
\numberwithin{equation}{section}
% Suite
\newcommand{\bs}[1]{\boldsymbol{#1}}
\newcommand{\dtemp}{\frac{d}{dt}}
\newcommand{\dpart}[2][]{\frac{\partial #1}{\partial #2}}
\newcommand{\inv}[1]{\frac{1}{#1}}
\newcommand{\card}[1]{\text{card}\,{#1}}

% Code
\usepackage{listings}

% Références
\usepackage{hyperref}

% Titre
\title{Formation \texttt{C++}}
\date{Septembre 2018}
\author[A.~\textsc{Tauveron}]{Aimery~\textsc{Tauveron-\,-Jalenques}}
\institute{\href{https://viarezo.fr/}{ViaRézo}}



\begin{document}

% Configuration pour listings
\lstset{frame=single, frameround=tttt, rulecolor=\color{blue}}
\lstset{breaklines=true, showstringspaces=false}
\lstset{language=[11]C++}

\frame{\titlepage}

% Logo VR
\addtobeamertemplate{frametitle}{}{
  \begin{tikzpicture}[remember picture,overlay]
    \node[anchor=north east,yshift=0pt] at (current page.north east) {\includegraphics[height=0.75cm]{VR_vertical.png}};
  \end{tikzpicture}
}

\section*{Introduction}
\label{sec:introduction}
\begin{frame}{Introduction}
  \begin{quotation}
    C++ est un langage de programmation compilé permettant la programmation sous de multiples paradigmes (comme la programmation procédurale, orientée objet ou générique). [\dots{}]

    Créé initialement par Bjarne Stroustrup dans les années 1980, le langage C++ est aujourd'hui normalisé par l'ISO. Sa première normalisation date de 1998, ensuite amendée par l'erratum technique de 2003. Une importante mise à jour a été ratifiée et publiée par l'ISO en septembre 2011 sous le nom de C++11. Depuis, des mises à jour sont publiées régulièrement : en 2014 puis en 2017.
  \end{quotation}

  \flushright{Extrait adapté de \href{https://fr.wikipedia.org/wiki/C\%2B\%2B}{Wikipédia}}
\end{frame}

\begin{frame}{Introduction}
  Ce n'est pas :
  \begin{itemize}
    \item une formation C (pas de \texttt{(int*)malloc(...)}),
    \item une formation POO,
    \item un tutoriel.
  \end{itemize}
\end{frame}


\frame{\tableofcontents}


\section{\texttt{C++} 101}
\label{sec:cpp101}
% -*- ispell-dictionary: "french" -*-

\begin{frame}[fragile]{Hello world!}
  \begin{lstlisting}
#include <iostream>

int main () {
  std::cout << "Hello world!\n";
  return 0;
}
  \end{lstlisting}
\end{frame}

\begin{frame}[fragile]{Compilation}
  Il faut compiler !
  \begin{itemize}
    \item GCC, Clang, Microsoft Visual C++
  \end{itemize}

  \begin{lstlisting}[language=bash]
g++ main.cpp
  \end{lstlisting}

  On peut exécuter :
  \begin{lstlisting}
./a.out
Hello world!
  \end{lstlisting}
\end{frame}

\begin{frame}[fragile]{Compilation}
  \begin{lstlisting}[language=bash]
g++ -O3 -Wall main.cpp -o hello
  \end{lstlisting}

  C'est mieux !
  \begin{lstlisting}
./hello
Hello world!
  \end{lstlisting}
\end{frame}

\begin{frame}[fragile]{Définition de variables}
  \begin{lstlisting}
#include <iostream>

int main () {
  int a;
  int b (0);
  int c = 0;
  std::cout << "a is: ";
  std::cin >> a;
  b = 2*a;
  std::cout << "a is: " << a << "\n";
  std::cout << "b is: " << b << "\n";
  std::cout << "2*a is: " << 2*a << "\n";
  return 0;
}
  \end{lstlisting}
\end{frame}

\begin{frame}{Les variables}
  Des types :
  \begin{itemize}
  \item numériques : \texttt{int} (\texttt{unsigned}, \texttt{long}, \texttt{short}), \texttt{float}, \texttt{double}
  \item non-numériques : \texttt{char}, \texttt{bool}
  \item \texttt{void}
  \item autres (objets\dots{})
  \end{itemize}

  Trois catégories de variables :
  \begin{itemize}
  \item variables : \textit{type} \textit{var}
  \item références : \textit{type} \texttt{\&}\textit{var}
  \item pointeurs : \textit{type} \texttt{*}\textit{var}
  \end{itemize}

  Un qualificateur : \texttt{const}
\end{frame}

\begin{frame}[fragile]{Les variables statiques}
  \begin{lstlisting}
int main () {
  int a (2);
  int b (a);
  int &c (a);
  int *d = a;
  std::cout << "a is: " << a << "\n";
  std::cout << "b is: " << b << "\n";
  std::cout << "c is: " << c << "\n";
  std::cout << "d is: " << d << "\n";
  return 0;
}
  \end{lstlisting}
\end{frame}

\begin{frame}[fragile]{Les variables statiques}
  \begin{lstlisting}
    int main () {
      int a (2);
      int b (a);
      int &c (a);
      int *d = &a; // address of a
      std::cout << "a is: " << a << "\n";
      std::cout << "b is: " << b << "\n";
      std::cout << "c is: " << c << "\n";
      std::cout << "d is: " << d << "\n";
      return 0;
    }
  \end{lstlisting}
\end{frame}

\begin{frame}[fragile]{Les variables statiques}
  \begin{lstlisting}
int main () {
  int a (2);
  int b (a);
  int &c (a);
  int *d = &a; // address of a
  std::cout << "a is: " << a << "\n";
  std::cout << "b is: " << b << "\n";
  std::cout << "c is: " << c << "\n";
  std::cout << "d is: " << *d << "\n"; /* the value pointed to by d */
  return 0;
}
  \end{lstlisting}
\end{frame}

\begin{frame}[fragile]{Les tableaux}
  \begin{lstlisting}
int a [8];
char w [256];
a[0] = 0;
a[1] = 4;
std::cout << a[0];
  \end{lstlisting}

  \begin{itemize}
  \item Hérités de C
  \item Indexés de 0 à $n-1$
  \item Des pointeurs $\implies$ pas de vérification des indices
  \end{itemize}

  \begin{lstlisting}
int *b = a; // Correct
a[11] = 32; // Undefined behaviour
  \end{lstlisting}  
\end{frame}

\begin{frame}[fragile]{Cycle de vie des variables statiques}
  \begin{itemize}
  \item Création à la déclaration
  \item Destruction à la fin du bloc \texttt{\{\dots\}}
  \end{itemize}

  \begin{lstlisting}
int main () {
  // maybe some code
  int a (0);
  unsigned int b;
  {
    double c (.5);
  }
  // c doesn't exist here
  return 0;
}
  \end{lstlisting}
\end{frame}

\begin{frame}[fragile]{Les variables dynamiques}
  Deux couples d'opérateurs : \texttt{new} et \texttt{delete}, \texttt{malloc} et \texttt{free}
  \begin{lstlisting}
int *a = new int(2);
delete a;
a = nullptr;
  \end{lstlisting}

  Pour les tableaux : \texttt{new[]} et \texttt{delete[]}
  \begin{lstlisting}
int *a = new int[4];
delete[] a;
a = nullptr;
  \end{lstlisting}
\end{frame}

\begin{frame}[fragile]{Les variables dynamiques}
  \begin{lstlisting}
int main () {
  int *a;
  {
    int *b = new int(4);
    a = b;
  }
  // b doesn't exist here
  // but a does
  std::cout << "a is: " << *a << "\n";
  delete a;
  a = nullptr;
  return 0;
}
  \end{lstlisting}
\end{frame}

\begin{frame}[fragile]{Les structures de contrôle}
  \begin{lstlisting}[escapechar=ù]
if ( ù\textit{condition}ù ) {
  // code
} else if ( ù\textit{condition}ù ) {
  // code
} else {
  // code
}
  \end{lstlisting}

  \begin{lstlisting}[escapechar=ù]
if ( ù\textit{condition}ù )
// single line of code
else
// single line of code
  \end{lstlisting}
\end{frame}

\begin{frame}[fragile]{Les structures de contrôle}
  \begin{lstlisting}[escapechar=ù]
while ( ù\textit{condition}ù ) {
  // code
}
  \end{lstlisting}
\end{frame}

\begin{frame}[fragile]{Les structures de contrôle}
  \begin{lstlisting}[escapechar=ù]
for ( ù\textit{initialisation}ù ; ù\textit{condition}ù ; ù\textit{increase}ù ) {
  // code
}
  \end{lstlisting}

  \begin{lstlisting}
for ( int i=0 ; i<10 ; ++i ) {
  // code
}
  \end{lstlisting}
  \begin{lstlisting}
for ( n=0, i=100 ; n!=i ; ++n, --i ) {
  // code
}
  \end{lstlisting}
\end{frame}

\begin{frame}[fragile]{Les structures de contrôle}
  \begin{lstlisting}[escapechar=ù]
switch ( ù\textit{expression}ù ) {
  case ù\textit{constant}ù:
  // code
  break;
  case ù\textit{constant}ù:
  // code
  break;
  default:
  //code
}
  \end{lstlisting}
\end{frame}

\begin{frame}[fragile]{Les fonctions}
  Définition :
  \begin{lstlisting}[escapechar=ù]
ù\textit{return\_type}ù ù\textit{function\_name}ù (ù\textit{arguments}ù) {
  // code
  return ù\textit{return\_value}ù;
}
  \end{lstlisting}

  Appel :
  \begin{lstlisting}[escapechar=ù]
ù\textit{variable}ù = ù\textit{function\_name}ù (ù\textit{arguments}ù);
  \end{lstlisting}

  Surcharge de fonctions possible
\end{frame}

\begin{frame}[fragile]{Les fonctions}
  Trois manières de passer un argument :
  \begin{itemize}
  \item par valeur
    \begin{lstlisting}
int f (int a) {...}
    \end{lstlisting}
  \item par référence
    \begin{lstlisting}
int f (int &a) {...}
    \end{lstlisting}
  \item par pointeur
    \begin{lstlisting}
int f (int *a) {...}
    \end{lstlisting}
  \end{itemize}
\end{frame}

\begin{frame}[fragile]{La fonction \texttt{main}}
  \begin{lstlisting}
int main (int argc, char **argv) {
  // code
  return 0;
}
  \end{lstlisting}
 ou bien
  \begin{lstlisting}
int main (int argc, char *argv[]) {
  // code
  return 0;
}
  \end{lstlisting}
\end{frame}

\begin{frame}{Des opérateurs}
  Dans le désordre :
  \begin{itemize}
  \item arithmétiques : +, -, *, /, \%
  \item de comparaison : ==, !=, <, >, <=, >=
  \item logiques : \&\&, ||, !
  \item binaires : \&, |, \^{}, \~{}
  \item mémoire : \&, *
  \end{itemize}

  En \texttt{C++}, le contexte est important !
\end{frame}

\begin{frame}[fragile]{Les \texttt{namespace}s}
  Pourquoi \texttt{std::cout} et non \texttt{cout} ? \texttt{int} et non \texttt{std::int} ?
  
  \textit{int} est dans le \texttt{namespace} global, \texttt{cout} dans le \texttt{namespace} \texttt{std}.
  
  \texttt{std} contient la bibliothèque standard.

  On peut utiliser :
  \begin{lstlisting}
using namespace std;

int main () {
  cout << "Hello world!\n";
  return 0;
}
  \end{lstlisting}
\end{frame}

\begin{frame}[fragile]{Les \texttt{namespace}s}
  \begin{lstlisting}
namespace foo {
  bool var (true);
}
namespace bar {
  double var (1.0);
}
int main () {
  std::cout << "foo::var is:" << foo::var << "\n";
  std::cout << "bar::var is:" << bar::var << "\n";
  return 0;
}
  \end{lstlisting}
\end{frame}



\section{Programmation orientée objet en \texttt{C++}}
\label{sec:poo}
% -*- ispell-dictionary: "french" -*-

\begin{frame}[fragile]{Structures et objets}
  Définir de nouveau types
  \begin{lstlisting}
    struct atom {
      int atomic; // Atomic number
      int mass; // Mass number
    };
  \end{lstlisting}

  Membres d'une fonction :
  \begin{itemize}
  \item variables (attributs),
  \item fonctions (méthodes)
  \end{itemize}
\end{frame}

\begin{frame}[fragile]{Structures et objets}
  \begin{lstlisting}
    atom carbon12;
    carbon12.atomic = 6;
    carbon12.mass = 12;

    atom *oxygen16 = new atom;
    oxygen16->atomic = 8;
    oxygen16->mass = 16;
  \end{lstlisting}
\end{frame}

\begin{frame}[fragile]{Structures et objets}
  \begin{lstlisting}
    class Atom {
    private:
      int atomic;
      int mass;
    public:
      Atom(int a, int m) {
        atomic = a; mass = m;
      }
    };

    atom carbon12 (6,12);
    carbon12.mass = 13; // Forbidden!
  \end{lstlisting}
\end{frame}

\begin{frame}{Sructures et objets}
  \textcolor{red}{Les attributs d'une classe sont toujours \textbf{privés}.}

  Les accès se font donc grâce aux méthodes de la classe.
\end{frame}

\begin{frame}[fragile]{Un objet minimaliste}
  \begin{lstlisting}
    class A {
    private:
      // attributes
      // private methods
    public:
      A(...) {...} // constructor
      ~A(...) {...} // destructor
      A(A &a) {...} // copy constructor
      // public methods
    };
  \end{lstlisting}
\end{frame}


\section{Le \texttt{C++} moderne}
\label{sec:cpp-moderne}
% -*- ispell-dictionary: "french" -*-

\begin{frame}[fragile]{Un peu d'histoire}
  \begin{description}
  \item[1972 :] première version de \texttt{C}
  \item[1983 :] première version de \texttt{C++} (alors \texttt{C with classes})
  \item[1989 :] premier standard pour \texttt{C} (\texttt{ANSI C})
  \item[1998 :] premier standard pour \texttt{C++}
  \item[2003 :] norme \texttt{C++03}
  \item[2011 :] norme \texttt{C++11} $\leftarrow$ début du \texttt{C++} moderne
  \item[2014 :] norme \texttt{C++14}
  \item[2017 :] norme \texttt{C++17}
  \end{description}
  \begin{lstlisting}[language=bash]
    --std=c++11
  \end{lstlisting}
\end{frame}

\begin{frame}[fragile]{De nouveaux tableaux}
  \begin{lstlisting}
std::array<int,5> a;
a[1] = 2;
size_t a_size = a.size();

a[6] = 4; // runtime error
  \end{lstlisting}
\end{frame}

\begin{frame}[fragile]{De nouveaux tableaux}
  \begin{lstlisting}
std::vector<int> a;
a.push_back(10);
a[0];

std::vector<int> b (5);
b.pop_back();
size_t b_size = b.size(); // now 4
  \end{lstlisting}
\end{frame}

\begin{frame}[fragile]{Une nouvelle boucle \texttt{for}}
  \begin{lstlisting}
std::array<int,5> a;
// code
for (int &it : a) {
  // code
}
  \end{lstlisting}
  \texttt{range for} pour tous les itérables
  \begin{itemize}
  \item \texttt{begin} et \texttt{end} sur l'objet
  \item \texttt{operator!=}, \texttt{operator++} et \texttt{operator*} sur l'itérateur
  \end{itemize}
\end{frame}

\begin{frame}[fragile]{Le type \texttt{auto}}
  \texttt{auto} : inférence de type
  
  \begin{lstlisting}[escapechar=ù]
for (auto &iterator : ù\textit{iterable}ù) {
  // code
}
  \end{lstlisting}
\end{frame}

\begin{frame}[fragile]{Parlons performance}
  Que se passe-t-il si j'exécute le code ci-dessous ?
  
  \begin{lstlisting}
std::string f () {
  // code
}

int main () {
  std::string x, y;
  // put some text in x and y
  std::string a(x);
  std::string b (x+y);
  std::string c (f());
  return 0;
}
  \end{lstlisting}
\end{frame}

\begin{frame}[fragile]{Parlons performance}
  Que se passe-t-il si j'exécute le code ci-dessous ?
  
  \begin{lstlisting}
std::vector<std::string> f () {
  // return a large vector of long strings
}

int main () {
  std::vector<std::string> a (f());
  return 0;
}
  \end{lstlisting}
\end{frame}

\begin{frame}{Parlons performance}
  Pourquoi faire une copie ?

  $\implies$ Indiquer au compilateur que l'objet affecté est détruit immédiatement après !
\end{frame}

\begin{frame}[fragile]{\textit{Move semantics}}
  \begin{itemize}
  \item Possible depuis \texttt{C++11}
  \item Utilisation automatique
  \item Utilisation manuelle avec \texttt{std::move}
  \end{itemize}
  
  \begin{lstlisting}[escapechar=ù]
ù\textit{new\_var}ù = std::move(ù\textit{old\_var}ù)
  \end{lstlisting}
\end{frame}

\begin{frame}[fragile]{\textit{Move semantics}}
  \begin{lstlisting}[escapechar=ù]
class A {
  // stuff
public:
  ~A ();
  A (const A &);
  A (A &&); // move constructor
  A& operator= (const A &);
  A& operator= (A &&); // move assignment operator
  // more stuff
};
  \end{lstlisting}
\end{frame}

\begin{frame}{\textit{Move semantics}}
  Autant que possible, utiliser les méthodes par défaut !
\end{frame}

\begin{frame}[fragile]{\textit{Move semantics} en pratique}
  \begin{lstlisting}
class blob {
private:
  size_t m_size;
  char *m_data;
public:
  blob (): m_size(0), m_data(nullptr) {}
  ~blob () { delete[] m_data; }
};
  \end{lstlisting}
\end{frame}

\begin{frame}[fragile]{\textit{Move semantics} en pratique}
  \begin{lstlisting}
blob (const blob& other): m_size(other.m_size) {
  // copy constructor
  if (m_size) {
    m_data = new char[m_size];
    memcpy(m_data, other.m_data, m_size);
  } else {
    m_data = nullptr;
  }
}
  \end{lstlisting}
\end{frame}

\begin{frame}[fragile]{\textit{Move semantics} en pratique}
  Avec \textit{copy-and-swap} :
  
  \begin{lstlisting}
friend void swap (blob& first, blob& second) {
  using std::swap;
  //
  swap(first.m_size, second.m_size);
  swap(first.m_data, second.m_data);
}

blob (blob&& other): blob() {
  // move constructor
  swap(*this, other);
}
  \end{lstlisting}
\end{frame}

\begin{frame}[fragile]{\textit{Move semantics} en pratique}
  Avec \textit{copy-and-swap} :
  
  \begin{lstlisting}
blob& operator= (blob other) {
  // assignment operator
  swap(*this, other);
  return *this;
}
  \end{lstlisting}

  Une seule définition nécessaire !
\end{frame}


\section{Introduction à la librairie standard}


\section*{Conclusion}
\label{sec:conclusion}
\begin{frame}{Ressources}
  \begin{itemize}
    \item \href{https://stackoverflow.com/}{Stack Overflow}
    \item \href{http://www.cplusplus.com/}{cplusplus.com}
  \end{itemize}
\end{frame}

\end{document}